% project-brief.tex — team-facing introduction to the SID injectivity experiment.
\documentclass[11pt]{article}

% --- page -------------------------------------------------------------------
\usepackage[a4paper,margin=1in]{geometry}
\usepackage{microtype}

% --- math -------------------------------------------------------------------
\usepackage{amsmath,amssymb,amsthm}
\DeclareMathOperator{\argmin}{arg\,min}

% --- tables / links ----------------------------------------------------------
\usepackage[hidelinks]{hyperref}
\usepackage{booktabs}   % \toprule \midrule \bottomrule
\usepackage{enumitem}   % tighter lists

\newcommand{\Keff}{K_{\mathrm{eff}}}
\newcommand{\Opi}{\mathcal{O}_\pi}
\newcommand{\Nitems}{N}

% --- title -------------------------------------------------------------------
\title{Semantic-ID Collisions in RQ-VAE Recommenders:\\ Experiment Brief and Metrics Guide}
\author{}
\date{November 2026}

\begin{document}
\maketitle

\section{The point}

We tokenize 4{,}783{,}154 item embeddings into Semantic IDs (SIDs) with a
Residual-Quantized VAE (RQ-VAE) using 3 codebook levels of 1024 codes each.
Nominal capacity is $1024^3 \approx 1.07 \times 10^9$ tuples, roughly 224 slots
per item, yet a large fraction of items share their full tuple. This experiment
asks \emph{why the capacity is not being realized}, by comparing training
schemes that differ in one axis only: how tightly the latent geometry and the
discrete assignment are coupled during training. Three arms are run
(`\texttt{rqvae}` baseline, `\texttt{warmup}`, `\texttt{drq}`), scored by the
same diagnostics, and reported here as raw numbers plus a guide to reading
them. The analysis itself is left open, and the logs and
`\texttt{*}\allowbreak diagnostics.json` files contain everything needed to do
it.

\section{Background}

\subsection{RQ-VAE and Semantic IDs}

An RQ-VAE maps an item embedding $x \in \mathbb{R}^{256}$ to a tuple of discrete
codes $(c_1,\dots,c_L)$, one per quantization level. The encoder produces a
continuous latent $z = E(x)$, then each level quantizes the residual left by the
previous levels:
\begin{equation}
  r_0 = z, \qquad
  c_\ell = \argmin_{i} \; \| r_{\ell-1} - e^{(\ell)}_i \|^2, \qquad
  r_\ell = r_{\ell-1} - e^{(\ell)}_{c_\ell},
  \label{eq:rq}
\end{equation}
where $e^{(\ell)}_i$ is the $i$-th entry of the level-$\ell$ codebook
$\mathcal{C}_\ell$, $|\mathcal{C}_\ell| = K$. The reconstructed latent is
$\hat{z} = \sum_{\ell=1}^{L} e^{(\ell)}_{c_\ell}$, decoded to $\hat{x} = D(\hat{z})$.
Training uses reconstruction plus the VQ objective with straight-through
gradients ($\mathrm{sg}$ is the stop-gradient operator):
\begin{align}
  \mathcal{L} &= \|x - \hat{x}\|^2
   + \sum_{\ell=1}^{L} \bigl[ \|\mathrm{sg}(r_{\ell-1}) - e^{(\ell)}_{c_\ell}\|^2
   + \beta \|r_{\ell-1} - \mathrm{sg}(e^{(\ell)}_{c_\ell})\|^2 \bigr],
  \label{eq:loss}\\
  z_{\mathrm{ste}} &= z + \mathrm{sg}(\hat{z} - z). \label{eq:ste}
\end{align}
Equation~\eqref{eq:ste} is how gradients bypass the non-differentiable
$\arg\min$ in \eqref{eq:rq}.

\emph{In plain terms:} the system gives every product a three-number barcode.
Products with similar content get similar barcodes: the first number is a broad
group label, the next two refine it. A generative recommender then learns to
predict that barcode, number by number, when it guesses what a user wants next.

\subsection{Why injectivity matters}

A generative recommender predicts the SID of the next item token by token, so
the SID map should be (nearly) one-to-one: items sharing a tuple become
indistinguishable to the model. The collision rate is the fraction of items that
do not have a private tuple. With $L=3$ and $K=1024$ we nominally have enough
slots, but the \emph{realized} capacity depends on how uniformly assignments
spread over tuples, and that depends on the latent geometry the encoder learns.
The tension: the reconstruction loss wants $z$ to concentrate (a good AE
compresses), while the quantizer wants assignments to spread over the codebook.
In joint training both fight through the same gradients.

\emph{In plain terms:} if two products share a barcode, the model cannot tell
them apart, so it cannot reliably serve either. The puzzle is the gap between
``there are enough barcodes'' (a billion possible) and ``they actually get
used'' (about a third of products share one). This project measures where that
gap comes from.

\subsection{The DRQ idea}

\textbf{DRQ} (Decoupled Residual Quantization, arXiv:2606.01844) splits the two
demands into sequential stages:
\begin{enumerate}[nosep]
  \item \emph{Stage 1, continuous VAE.} Train $E, D$ with reconstruction only
  (optionally plus a contrastive term $\lambda \mathcal{L}_{\mathrm{InfoNCE}}$;
  $\lambda = 0$ in our run). No codebook, no straight-through gradient. The
  latent geometry settles without quantizer pressure.
  \item \emph{Stage 2, post-hoc hierarchical $k$-means.} Freeze the latents and
  fit each level's codebook by dataset-level $k$-means on the residuals of that
  level. Centroids are means over all points assigned to them, so every point
  influences every centroid at every iteration, unlike STE training where only
  the argmin-selected code receives a gradient.
\end{enumerate}
The third arm, \textbf{warmup}, is a control between the two: it runs the stage-1
AE for 5 epochs alone, then hands off to the \emph{identical} joint RQ-VAE
training as the baseline. It isolates the effect of ``unperturbed latent first''
from the effect of ``$k$-means assignment instead of learned codebooks''.

\emph{In plain terms:} instead of teaching the map-maker and the postcode
system at the same time and letting them interfere, DRQ first lets the network
draw the map of all products however it likes, then freezes it and stamps
postcodes on the finished map with a classic grouping routine. The warmup arm
splits the difference: clean map first, but postcodes are still learned.

\section{Experimental setup}

Both the arms and the scoring are implemented in
`\texttt{train\_drq\_vs\_rqvae.py}` (the archived copies in each run folder are
the exact scripts that produced those numbers). Shared configuration:

\begin{itemize}[nosep]
  \item Data: 4{,}783{,}154 items, 256-dim embeddings (parquet);
        train 4{,}304{,}839, validation 478{,}315 (10\% split, seed 42).
  \item Architecture: encoder $256 \to \cdots \to 64$, decoder symmetric,
        $L=3$ levels, $K=1024$ codes per level, latent dim 64.
  \item Joint training (arms `\texttt{rqvae}` and `\texttt{warmup}`): 25 epochs,
        batch 4096, Adam lr $5\times10^{-4}$, 2000-step linear warmup, cosine
        decay to 1\% floor, $\beta = 0.1$, trunc-normal codebook init,
        dead-code revival enabled (revive codes used $<10\%$ of mean usage,
        checked every 500 steps, until 60\% of the step budget).
        The `\texttt{warmup}` arm prepends 5 reconstruction-only epochs and
        restarts the cosine schedule with a 200-step ramp at the handoff.
  \item DRQ arm: stage 1 uses the same optimizer schedule for 25 epochs with
        $\lambda = 0$; stage 2 fits FAISS GPU $k$-means (50 iterations,
        3 restarts) per level over the full frozen latent pool, then assigns
        all items level by level through the residual cascade
        \eqref{eq:rq}.
  \item Diagnostics are computed identically for every arm on the full 4.78M
        index, after training. SIDs themselves never depend on any diagnostic
        parameter (in particular the kernel width $\sigma$ in
        section~\ref{sec:overlap} is a post-hoc metric choice only).
\end{itemize}

\section{The metrics}

This is the section to read before opening any diagnostics file. Metrics come in
three tiers: per-level codebook health, tuple-level injectivity, and
reconstruction. The full pipeline claim of DRQ is that better per-level health
and better geometry reduce collisions; checking \emph{which} tier moved is how
you test that claim.

\emph{In plain terms:} three questions, in order. \emph{Are all the bins being
used?} (per-level metrics: a bin nobody puts anything in, or a bin holding a
third of everything, is wasted capacity). \emph{Are the bins genuinely
different, or do some sit on top of each other?} (geometry metrics). \emph{Do
products end up with private barcodes in the end?} (tuple metrics: the business
number). Reconstruction asks whether the barcodes still describe the product
accurately, the quality tax of the whole exercise.

\subsection{Per-level (marginal) metrics}

For level $\ell$, let $\pi_i$ be the fraction of the 4.78M items assigned to
code $i$, over $K = 1024$ codes.

\begin{itemize}[nosep]
  \item \textbf{used codes / dead \%}: $|\{i : \pi_i > 0\}|$. Necessary but not
  sufficient: a level can use every code and still be useless jointly.
  \item \textbf{entropy / perplexity}: $H_\ell = -\sum_i \pi_i \log \pi_i$ (nats),
  $\mathrm{perp}_\ell = e^{H_\ell}$. Perplexity reads as ``effective number of
  balanced codes''; 1024 is the perfectly uniform ceiling.
  \item \textbf{Gini} and \textbf{max\_freq}: inequality of the usage distribution;
  0 and $1/K$ are uniform.
  \item \textbf{distribution floor}: $\sum_i \pi_i^2$, the self-collision mass of
  the level's own assignment distribution if that level were the whole ID.
  \item \textbf{participation ratio}: $1 / \sum_i \pi_i^2$, the inverse of the
  floor. Pure usage-balance view of capacity; ignores centroid positions.
\end{itemize}

\emph{In plain terms:} this tier answers ``is every level pulling its weight?''.
Healthy values here look like a warehouse where every aisle and every shelf is
occupied roughly equally. What these metrics \emph{cannot} see is whether the
same handful of choices repeat across levels; that is what the joint tier is
for.

\subsection{Overlap rate and effective codes $\Keff$ (geometry)}
\label{sec:overlap}

Two codes that are used but sit on top of each other are one code in disguise.
DRQ measures this with the \emph{expected overlap rate}, the probability that an
item and an independent second draw from the same level land within each other's
Gaussian kernel (bandwidth $\sigma$):
\begin{equation}
  \Opi^{(\ell)}(\sigma) \;=\; \sum_{i}\sum_{j} \pi_i \pi_j\,
    \exp\!\Bigl(-\frac{\|c_i - c_j\|^2}{4\sigma^2}\Bigr),
  \qquad
  \Keff^{(\ell)} \;=\; \frac{1}{\Opi^{(\ell)}} .
  \label{eq:opiq}
\end{equation}
It decomposes exactly into the two failure modes:
\begin{equation}
  \underbrace{\Opi = \sum_i \pi_i^2}_{\text{distribution floor}}
  \;+\;
  \underbrace{\sum_{i\neq j} \pi_i\pi_j \exp\!\bigl(-\|c_i-c_j\|^2/4\sigma^2\bigr)}
  _{\text{geometry cross-overlap}} .
  \label{eq:split}
\end{equation}
The floor term is usage skew (fixable by any balance mechanism); the cross term
is codebook geometry (centroids crowding the manifold), which is what DRQ
stage-2 is argued to fix. In the JSON, \texttt{distribution\_floor} and
\texttt{geometry\_cross\_overlap} are reported separately.

$\sigma$ is chosen per level automatically: 0.1 $\times$ the median
nearest-neighbour distance among used centroids
(\texttt{kernel\_sigma} in the JSON). The diagnostics also include a
\texttt{sigma\_scan} over $\sigma$ at $\times0.25, 0.5, 1, 2, 4$ of that auto
value. Read the scan as ``how close do centroids have to be before they start
counting as duplicates'': $\Keff$ should stay near the participation ratio at
small-to-moderate $\sigma$ and collapse as $\sigma$ grows past the typical
centroid spacing. A codebook with real geometric crowding collapses early in
the scan; a spread one holds $\Keff \approx$ PR until the kernel covers
half the space. Comparing the two arms' scan curves is the cleanest geometry
test in the data.

\emph{In plain terms:} this asks whether distinct bin numbers really mark
distinct regions, or whether two bins cover the same territory and split it
arbitrarily. The scan varies the notion of ``close'' (from very strict to quite
forgiving) and watches how many bins remain distinct. Bins that merge under a
lax definition were nearly duplicates all along.

\subsection{Tuple-level (joint) metrics}

These measure what the recommender actually sees, the map item $\to$ $(c_1,c_2,c_3)$.

\begin{itemize}[nosep]
  \item \textbf{collision\_rate} $= 1 - U/N$, with $U$ the number of unique
  full tuples and $N$ the item count. The headline metric.
  \item \textbf{n\_unique\_tuples}, \textbf{max\_collision\_bucket} (largest
  group of items sharing one tuple; a tail metric), \textbf{mean\_collision\_bucket},
  \textbf{items\_in\_collision} (fraction of items that are in any shared bucket).
  \item \textbf{joint\_entropy / joint\_perplexity}: entropy of the empirical
  distribution over tuples, and $e^H$. Joint perplexity is the size of the
  effectively-used tuple space.
  \item \textbf{effective\_cells}: the same relation inverted. Given $U/N$, find
  $x$ with $(1 - e^{-x})/x = U/N$ and report $M_{\mathrm{eff}} = N/x$. It
  answers ``how many uniformly-sized tuple cells would produce this many
  collisions?''.
  \item \textbf{redundancy\_bits} $= \bigl(\sum_\ell H_\ell\bigr) - H_{\mathrm{joint}}$
  (converted to bits). How much of the nominal $3 \times \log_2 1024 = 30$ bits
  of tuple space is wasted because levels repeat each other's information.
  A perfectly independent cascade keeps redundancy low; strong residual
  correlation between levels raises it.
  \item \textbf{prefix\_uniqueness}[d]: unique fraction of length-$d$ prefixes.
  $d=1$ is bounded by $K/N$; the $d=2 \to d=3$ jump shows where collisions are
  decided. If depth-3 prefix uniqueness $\approx$ collision-implied uniqueness,
  the last level is doing the de-duplicating work (or failing at it).
\end{itemize}

\emph{In plain terms:} this is the tier the business actually buys: what
fraction of the catalog has a private barcode (collision rate), how bad the
worst pile-up is (max bucket), and how large the effectively-used barcode space
is (joint perplexity, effective cells). Redundancy measures how many digits are
spent saying things an earlier digit already said.

\subsection{Reconstruction metrics}

\begin{itemize}[nosep]
  \item \textbf{val recon (continuous)}: $\|x - D(E(x))\|^2$ on held-out items,
  the AE quality ceiling.
  \item \textbf{val recon (quantized)}: $\|x - D(\hat z)\|^2$ through the codes.
  \item \textbf{quantization cost}: quantized $-$ continuous, the price paid for
  discretization. This is the fairness axis: a scheme can improve collisions
  while paying with reconstruction. All terms are summed over the 256 dims and
  averaged over the batch.
\end{itemize}

\emph{In plain terms:} does the barcode system still describe products well?
The quantization cost is the price of discretization: schemes that win on
collisions can quietly pay with worse descriptions, so read this tier next to
the collision tier, never alone.

\subsection{How to think about the numbers}

A sizing sanity check before looking at any arm: if items landed
\emph{uniformly at random} in $M$ tuple cells, the expected unique fraction is
\begin{equation}
  \frac{U}{N} \;=\; \frac{1 - e^{-x}}{x}, \qquad x = \frac{N}{M}.
  \label{eq:birthday}
\end{equation}
With nominal $M = 1024^3$, $x \approx 0.0045$ and \eqref{eq:birthday} predicts
$U/N \approx 0.998$: uniform use of the nominal space would be essentially
collision-free. Any measured uniqueness far below that comes from either
uneven per-level marginals (bounded by $\prod_\ell \Keff^{(\ell)}$ at small
$\sigma$) or joint concentration (redundancy, prefix curves), so the three
tiers of metrics above are exactly the three candidate explanations. Questions
worth answering from the logs: which tier differs between arms; does the
$\Keff$ ladder $\mathrm{rqvae} < \mathrm{warmup} < \mathrm{drq}$ line up with
the collision ladder; and what does \eqref{eq:birthday} say about each arm's
effective cells.

\section{Results}

Raw numbers only. All three arms share the configuration of section 3 and are
scored by the same diagnostics, including the multi-$\sigma$ scan; the baseline
rerun reproduced its earlier numbers exactly (same seed).

\subsection{Tuple-level}

\begin{table}[h]
\centering
\begin{tabular}{lrrr}
\toprule
metric & rqvae & warmup & drq \\
\midrule
collision\_rate           & 0.28330 & 0.28635 & 0.25726 \\
unique tuples             & 3{,}428{,}074 & 3{,}413{,}491 & 3{,}552{,}635 \\
max collision bucket      & 564    & 317    & 196    \\
mean collision bucket     & 1.395  & 1.401  & 1.346  \\
items in collision        & 1{,}942{,}728 & 1{,}958{,}503 & 1{,}822{,}616 \\
joint perplexity          & 2{,}621{,}981 & 2{,}600{,}081 & 2{,}873{,}590 \\
effective cells           & 6{,}758{,}007 & 6{,}666{,}973 & 7{,}622{,}188 \\
redundancy (bits)         & 8.315  & 8.341  & 8.272  \\
prefix uniqueness $d{=}2$ & 0.1161 & 0.1145 & 0.1166 \\
prefix uniqueness $d{=}3$ & 0.7167 & 0.7136 & 0.7427 \\
\bottomrule
\end{tabular}
\caption{Tuple-level diagnostics, $N = 4{,}783{,}154$, $L=3$, $K=1024$, seed 42.}
\end{table}

\subsection{Per-level}

\begin{table}[h]
\centering
\begin{tabular}{llrrrr}
\toprule
arm & level & perp & Gini & $\Keff$ (auto $\sigma$) & dist.\ floor \\
\midrule
rqvae  & 0 & 951 & 0.212 & 887.5 & 0.00113 \\
rqvae  & 1 & 912 & 0.268 & 821.4 & 0.00122 \\
rqvae  & 2 & 962 & 0.193 & 902.6 & 0.00111 \\
\midrule
warmup & 0 & 957 & 0.203 & 899.2 & 0.00111 \\
warmup & 1 & 911 & 0.268 & 820.9 & 0.00122 \\
warmup & 2 & 967 & 0.186 & 910.6 & 0.00110 \\
\midrule
drq    & 0 & 975 & 0.173 & 930.3 & 0.00108 \\
drq    & 1 & 950 & 0.213 & 887.8 & 0.00113 \\
drq    & 2 & 958 & 0.195 & 910.2 & 0.00110 \\
\bottomrule
\end{tabular}
\caption{Marginal health per level; all levels use 1024/1024 codes in all arms.}
\end{table}

\subsection{Reconstruction}

\begin{table}[h]
\centering
\begin{tabular}{lrrr}
\toprule
& rqvae & warmup & drq \\
\midrule
val recon (continuous) & 0.01794 & 0.01613 & 0.01305 \\
val recon (quantized)  & 0.08262 & 0.08193 & 0.10130 \\
quantization cost      & 0.0647  & 0.0658  & 0.0883 \\
\bottomrule
\end{tabular}
\caption{Validation reconstruction (sum over 256 dims, mean over batch). These
three values are printed in `\texttt{run.log}`; they are not part of the JSON.}
\end{table}

\subsection{Sigma scan ($\Keff$ under widening kernels)}

\begin{table}[h]
\centering
\begin{tabular}{llrrrrr}
\toprule
arm & level & PR & $\times0.25$ & $\times0.5$ & $\times1$ & $\times4$ \\
\midrule
rqvae  & 0 & 888 & 888 & 888 & 888 & 279 \\
rqvae  & 1 & 821 & 821 & 821 & 821 & 37 \\
rqvae  & 2 & 903 & 903 & 903 & 903 & 22 \\
\midrule
warmup & 0 & 899 & 899 & 899 & 899 & 279 \\
warmup & 1 & 821 & 821 & 821 & 821 & 36 \\
warmup & 2 & 911 & 911 & 911 & 911 & 22 \\
\midrule
drq    & 0 & 930 & 931 & 930 & 930 & 218 \\
drq    & 1 & 888 & 888 & 888 & 888 & 51 \\
drq    & 2 & 910 & 910 & 910 & 910 & 26 \\
\bottomrule
\end{tabular}
\caption{$\Keff^{(\ell)}(\sigma)$ at $\sigma$ = mult $\times$ auto value
(auto $\sigma\approx$ 0.06 / 0.03 / 0.02 for levels 0/1/2 in all arms).
PR column is the participation ratio, the usage-only ceiling that
$\Keff\le$PR can never exceed; the $\times2$ column is omitted for space, its
values: rqvae 872/703/782, warmup 883/699/800,
drq 910/826/846. At the auto value the geometry cross-overlap term is
$\le 2.5\times10^{-9}$ for every arm and level.}
\end{table}

The `\texttt{drq}` column shows the same collision-rate reduction direction
reported for DRQ-VAE in its paper (0.229 $\to$ 0.075 on the paper's corpus and
codebook sizes; ours is 0.283 $\to$ 0.257 on 4.78M items at 3$\times$1024).
Whether that is the whole story, or which tier of the metrics explains it, is
the open question for the group.

\section{Files you will receive}

\begin{itemize}[nosep]
  \item `\texttt{logs/rqvae\textless jobid\textgreater.out}`: full stdout,
  including the per-epoch table and the handoff line for `\texttt{warmup}`.
  \item `\texttt{arm\_\{rqvae,warmup,drq\}\_diagnostics.json}`: the scored tiers
  of section 4. Keys: `\texttt{tuples}` (joint), `\texttt{stage}` (per level,
  with `\texttt{sigma\_scan}`), `\texttt{prefix\_uniqueness}`,
  `\texttt{redundancy\_bits}`. Reconstruction values are in `\texttt{run.log}`.
  \item `\texttt{comparison*.json}`: pairwise tables built by the script when
  two arm JSONs share a run folder.
  \item `\texttt{arm\_*\_index.json}`: item\_id $\to$ SID tuples, usable for
  any custom collision analysis.
  \item `\texttt{train\_drq\_vs\_rqvae.py}`: exact code that produced the run.
\end{itemize}

\section{References}

\begin{enumerate}[nosep,leftmargin=*]
  \item DRQ: \emph{Rethinking Semantic ID Construction for Generative
        Recommendation} (arXiv:2606.01844). Basis of the `\texttt{drq}' arm and
        the overlap metrics.
  \item TIGER: \emph{Recommender Systems with Generative Retrieval}
        (arXiv:2305.05065). The RQ-VAE SID recipe being improved on.
  \item GRID handbook (arXiv:2507.22224). Reports RQ-VAE's weakness against
        simpler alternatives; motivates this study.
  \item QuaSID (arXiv:2603.00632) and FORGE (arXiv:2509.20904). Post-hoc
        collision-handling alternatives.
\end{enumerate}

\end{document}
